\subsection{Why grouping regularizes the readout}
\label{sec:mechanism}

Full averaging gives the readout one clean view of each time step. Individual
shots give it many views, but every view is noisy. Intermediate groups provide
several partially denoised versions of the same target. Differences among
those versions reveal which feature changes are stable across shots and which
are measurement noise. Training on all grouped versions discourages the
readout from relying on the unstable changes.

For a linear readout with squared loss, this intuition can be written as an
exact identity. Let $\bar{x}_t$ be the mean of the $G$ grouped
views and let
$C_{k,t}=G^{-1}\sum_g(x_t^{(g)}-\bar{x}_t)(x_t^{(g)}-\bar{x}_t)^\top$.
This quantity summarizes how the measured features fluctuate together around
their mean. For intercept $b$ and weights $w$,
\begin{equation}
 \frac{1}{G}\sum_g\!\left(y_t-b-w^\top x_t^{(g)}\right)^2
 =\left(y_t-b-w^\top\bar{x}_t\right)^2+w^\top C_{k,t}w .
 \label{eq:covariance}
\end{equation}
The first term on the right is the EV loss. The second adds a larger cost when
the readout relies on combinations of features that vary strongly across the
grouped views. Grouped training is therefore exactly EV training plus a
penalty determined by the observed shot fluctuations. This identity requires
neither Gaussian noise nor independence across time.

Cov-EV applies the same fluctuation-based penalty directly to EV and chooses
its strength separately using the same earlier data as Grouped. If both methods
perform similarly, the benefit of grouping is explained by this regularization
rather than by additional physical information. We also compute both forms at
fixed $k$ and confirm identical predictions as an implementation check.
The leaky integrator and overlapping lags make adjacent feature sets
dependent, so uncertainty is computed over trajectories or measured windows,
not over the many grouped views. The full derivation and statistical details are
provided in the supplement.
